\section{Gregorian calendar continuity and exact phase carry}
\label{sec:calendar-phase}

One January-to-December sequence does not cover the Gregorian leap rules. The
calendar proceeds into the next year; a wrapped address phase does not restart
the calendar. The new verification enumerates every date from 1 January 2000
through 31 December 2399 and includes the following 1 January 2400 boundary.
This is a calendar and address check, separate from satellite forecast accuracy.

\subsection{Independent date arithmetic}
For year $Y$, month $M$ and day $D$, the arithmetic oracle is
\begin{align*}
a&=\left\lfloor(14-M)/12\right\rfloor,\qquad
y=Y+4800-a,\qquad m=M+12a-3,\\
J&=D+\left\lfloor(153m+2)/5\right\rfloor+365y\\
 &\quad+\left\lfloor y/4\right\rfloor-\left\lfloor y/100\right\rfloor
 +\left\lfloor y/400\right\rfloor-32045,\\
T&=86400\,[J-J(1980,1,6)].
\end{align*}
Here $J$ is the astronomical-noon Julian day number associated with a date;
integer date differences produce the corresponding midnight tick $T$.
Expected dates use explicit month enumeration, checked against both this
March-based formula and an independent January-based ordinal sum. The expected
value oracle does not use the runtime's \texttt{datetime} implementation.

A year is leap when divisible by four, except centuries not divisible by 400.
The complete cycle therefore has
\[
97=100-4+1\quad\hbox{leap days},\qquad
400\cdot365+97=146097=7\cdot20871\quad\hbox{days}.
\]
The month/leap pattern and weekdays repeat after 400 years. The year number and
absolute time continue. For each enumerated date, the arithmetic difference
to the same date 400 years later is also checked; this does not add another
400 years of formatter or physical-orbit coverage.

\subsection{The phase wraps; absolute time continues}
For integer reference $T_0$ and period $P=2^{14}$, the literal address helper uses
\[
(w,r)=\operatorname{divmod}(T-T_0,P),\qquad
T=T_0+Pw+r,\quad 0\le r<P,\qquad f=r/P.
\]
Arithmetic right shift and a 14-bit mask form the independent test oracle.
The fraction $f$ is exactly representable in binary64. All phase values are
tested at three signed winding offsets, with additional signed tick/reference
cases reaching $10^{100}+16385$. This arbitrary-size claim concerns the integer
helper, not binary64 timestamp input.

The contiguous and Morton key layouts store the remainder $r$ in their time
field. A key may repeat after $P$ ticks while winding $w$ increases. Tests check
both layouts against independently assembled bit strings and retain absolute
tick, winding and remainder in every daily chart record. At a one-second tick
period, the phase repeats after 16,384 seconds; it is not a calendar month or year.

\subsection{Recorded result and limits}
\begin{center}\begin{tabular}{@{}lr@{}}\toprule
Check coverage & Count\\\midrule
Dates in the complete cycle & 146,097\\
Daily midnights including the next boundary & 146,098\\
Month starts / year starts & 4,800 / 400\\
Leap days & 97\\
Daily GPST formatter comparisons & 438,293\\
Additional century/epoch formatter comparisons & 22\\
Phase wraps / adjacent tested ticks & 770,433 / 2,311,299\\
Complete phase cases at three signed offsets & 49,152\\
Large signed tick/reference cases & 210\\
Complete daily chart/key records & 146,098\\\midrule
Exact comparisons / failures & 9,519,315 / 0\\\bottomrule
\end{tabular}\end{center}

Seven focused regression tests also pass, including checks that deliberately
broken calendar-reset and lost-winding substitutes are detected. Explicit
century cases cover common years 1900, 2100, 2200 and 2300, and leap years 2000
and 2400. The 2400 leap-day cases are separate checks outside the enumerated
cycle. Daily checks exercise midnight, an exact 0.125-second fraction, and
the preceding date's last whole second.

The machine-readable result is
\path{results/calendar_phase_36110/summary.json}; it records the checked source
hashes and all eight unchanged seed hashes. Reproduce it with
\texttt{python tools/validate\_calendar\_phase.py}.
No runtime equations, floating-point propagation or force-model values changed.

\begin{panel}
The formatter labels uniform 86,400-second days as GPST. This audit does not
validate a UTC leap-second table, predict future UTC offsets, certify
sub-millisecond display accuracy, or expand any seed's physical timestamp
domain. The calendar check supplies no 400-year orbit forecast or additional
satellite-position accuracy guarantee.
\end{panel}
